\documentclass[11pt,a4paper]{article}

% ---------------------------------------------------------------------
%  Paquetes (todos vienen en TeX Live / MiKTeX / Overleaf)
% ---------------------------------------------------------------------
\usepackage[utf8]{inputenc}
\usepackage[T1]{fontenc}
\usepackage{lmodern}
\usepackage[margin=2.2cm]{geometry}
\usepackage{amsmath}
\usepackage{xfp}            % cálculo de la nota sugerida
\usepackage{array,booktabs,longtable}
\usepackage[table]{xcolor}
\usepackage{graphicx}
\usepackage{hyperref}
\hypersetup{colorlinks=true, linkcolor=black, urlcolor=blue}
\setlength{\parindent}{0pt}
\setlength{\parskip}{0.5em}
\renewcommand{\contentsname}{Contenido}

% =====================================================================
%  macros.tex — comandos del informe GA-EXP-02 (no hace falta editarlo)
% =====================================================================

% Texto pendiente: aparece en ROJO en el PDF para que no se olvide nada.
\newcommand{\completar}[1]{\textcolor{red}{\textbf{[COMPLETAR: #1]}}}
\newcommand{\faltanivel}{\textcolor{red}{\textbf{?}}}

% ---------- Identificación del equipo evaluado ----------
% \identificacion{Equipo}{Tema}{Fecha}{Hora inicio}{Hora fin}
\newcommand{\identificacion}[5]{%
  \subsection*{Identificación}
  \begin{tabular}{@{}p{0.28\textwidth}p{0.66\textwidth}@{}}
    \toprule
    Equipo evaluado & #1 \\
    Tema            & #2 \\
    Fecha           & #3 \\
    Hora de inicio  & #4 \\
    Hora de fin     & #5 \\
    \bottomrule
  \end{tabular}
}

% ---------- Tabla de dimensiones ----------
\newenvironment{dimensiones}{%
  \subsection*{Evaluación por dimensión}
  \renewcommand{\arraystretch}{1.25}
  \begin{longtable}{@{}>{\raggedright\arraybackslash}p{0.24\textwidth}>{\centering\arraybackslash}p{0.07\textwidth}>{\raggedright\arraybackslash}p{0.63\textwidth}@{}}
    \toprule
    \textbf{Dimensión (peso)} & \textbf{Nivel} & \textbf{Evidencia observable} \\
    \midrule
    \endhead
}{%
    \bottomrule
  \end{longtable}
}
% \dimension{Nombre (peso)}{Nivel 0-4}{Evidencia}
\newcommand{\dimension}[3]{#1 & #2 & #3 \\ \midrule}

% ---------- Nota sugerida con la cuenta escrita ----------
% \notasugerida{D1}{D2}{D3}{D4}{D5}{D6}   (niveles 0-4)
% Fórmula de la guía: (10*D1 + 20*D2 + 15*D3 + 20*D4 + 25*D5 + 10*D6) / 4 * 0.1
\newcommand{\notasugerida}[6]{%
  \subsection*{Nota sugerida}
  \edef\NSsuma{\fpeval{10*(#1)+20*(#2)+15*(#3)+20*(#4)+25*(#5)+10*(#6)}}%
  \edef\NSdiv{\fpeval{\NSsuma/4}}%
  \edef\NSnota{\fpeval{\NSdiv*0.1}}%
  Fórmula: $\dfrac{10D_1+20D_2+15D_3+20D_4+25D_5+10D_6}{4}\times 0.1$
  \[
    \frac{10(#1)+20(#2)+15(#3)+20(#4)+25(#5)+10(#6)}{4}\times 0.1
    = \frac{\NSsuma}{4}\times 0.1 = \NSdiv\times 0.1 = \NSnota
    \approx \mathbf{\fpeval{round(\NSnota,1)}}
  \]
}

% ---------- Retroalimentación ----------
% \retroalimentacion{Fortaleza}{Mejora}
\newcommand{\retroalimentacion}[2]{%
  \subsection*{Fortaleza}
  #1
  \subsection*{Mejora accionable}
  #2
}


% ---------------------------------------------------------------------
%  Datos de la carátula
% ---------------------------------------------------------------------
\newcommand{\estudiante}{Mera Jumbo Allison Michelle}
\newcommand{\fechasesion}{26 de noviembre de 2026}

\begin{document}

% ============================ CARÁTULA ===============================
\pagenumbering{gobble}
\begin{titlepage}
  \centering
  {\Large\bfseries Universidad Técnica Estatal de Quevedo\par}
  \vspace{0.3cm}
  {\large Facultad de Ciencias de la Computación\par}
  {\large Carrera de Ingeniería de Software\par}
  \vspace{2cm}
  {\large ISR-701 Aplicaciones Distribuidas\par}
  \vspace{0.5cm}
  {\Huge\bfseries GA-EXP-02\par}
  \vspace{0.3cm}
  {\LARGE Coevaluación de exposiciones\par}
  \vspace{2cm}
  {\large\textbf{Estudiante:} \estudiante\par}
  \vspace{0.3cm}
  {\large\textbf{Docente:} Gleiston Cicerón Guerrero Ulloa, PhD\par}
  \vspace{0.3cm}
  {\large\textbf{Fecha de la sesión:} \fechasesion\par}
  \vfill
  {\large Séptimo semestre --- SPA 2026--2027\par}
\end{titlepage}

\pagenumbering{arabic}
\tableofcontents
\clearpage

% ======================= EQUIPOS OBSERVADOS ==========================
% Un archivo por equipo (todos menos el propio). Para agregar otro:
% copia secciones/_plantilla_equipo.tex como equipo_NN.tex y añade su \input.
% =====================================================================
%  EQUIPO 1 (GGHL) — T7 Computación paralela — 26/11/2026
%  Integrantes observados: Diego, Roger, Lombeida.
%  Las diapositivas citadas (1/5 a 5/5) son las de la parte de Diego.
%  PENDIENTE: poner los niveles (?) y los mismos números en \notasugerida.
% =====================================================================

\section{Grupo 1 (GGHL) --- T7 Computación paralela}

\identificacion
{Grupo 1 (GGHL): Diego, Roger y Lombeida}
{T7 --- Computación paralela (Apache Spark)}
{26/11/2026}
{08:40}
{09:10}

\begin{dimensiones}
	\dimension{D1 Cumplimiento de la guía (10)}
	{3}
	{Conceptos clave: presentados por Diego (diap.~1--2). Demostración en vivo:
		Diego ejecutó la práctica RTV con Spark \texttt{local[*]} (diap.~3);
		Lombeida presentó una práctica pequeña; Roger no presentó práctica.
		Comparación: solo entre 2 y 4 particiones de la misma práctica (diap.~3).
		Recomendación: no se observó una recomendación explícita.}
	
	\dimension{D2 Fundamentación conceptual (20)}
	{3}
	{Diego explicó: definición de computación paralela y su objetivo
		(diap.~1); diferencia entre ejecución secuencial, paralela y distribuida con
		el ejemplo de vehículos (diap.~1); arquitectura de Spark con driver, gestor
		del clúster, nodos y ejecutores (diap.~2); tareas en hilos y núcleos
		(diap.~2); particiones y las operaciones \texttt{parallelize()},
		\texttt{map()} y \texttt{collect()} (diap.~3). No se detectaron errores
		graves de la guía. La explicación de Roger quedó incompleta.}
	
	\dimension{D3 Desarrollo técnico (15)}
	{2}
	{Diego explicó el código y lo ejecutó: Spark en modo \texttt{local[*]}
		en una sola computadora (diap.~3), es decir, no en procesos o equipos
		distintos, y sin caso de fallo. Resultado: 2 particiones, 2/2 tareas,
		1,8195\,s; 4 particiones, 4/4 tareas, 3,4206\,s (dos mediciones por
		configuración). Se comparó una sola variable (número de particiones) con
		un criterio medible (tiempo), no dos alternativas tecnológicas. Roger no
		hizo demostración de código; la práctica de Lombeida fue pequeña y su
		explicación del código, escasa.}
	
	\dimension{D4 Análisis y recomendación (20)}
	{3}
	{Se reconoció una limitación: en la prueba, aumentar las particiones
		elevó el tiempo, y en un clúster también influyen el reparto y la
		ubicación de los datos (diap.~3). No se indicó cuándo usar y cuándo no
		usar la computación paralela ni se ligaron razones a un caso concreto.}
	
	\dimension{D5 Comunicación (25)}
	{3}
	{Duración: 30\,min (08:40--09:10), sin tiempo asignado para comparar con
		$T \pm 2$\,min. Expusieron los tres integrantes. Diego explicó su parte con
		claridad. A Lombeida no se le entendió bien por su entonación. No hubo
		preguntas, así que no se observó respuesta a una pregunta técnica.}
	
	\dimension{D6 Fuentes (10)}
	{2}
	{En las diapositivas de Diego, cada dato tiene su cita visible (diap.~1--3),
		incluida la tabla de tiempos (``registro de la práctica''). Las referencias
		están en formato IEEE con DOI (diap.~4--5): fuentes indexadas (Springer,
		\emph{IEEE Transactions}, \emph{The Computer Journal}) y documentación
     oficial de Apache Spark. El dato del 80\,\% de Fortune 500 (diap.~1)
     proviene de la página del propio proyecto. No se verificó en las diapositivas de Roger y Lombeida.}
\end{dimensiones}

% Pon aquí los mismos niveles de la tabla, en orden D1..D6 (solo números):
\notasugerida{3}{3}{2}{3}{3}{2}

\retroalimentacion
  {Por parte de Diego buen dominio del tema y de explicación}
  {Por parte de los compañeros Roger y Lombeida se recomienda profundizar en la explicación práctica.}

\clearpage

%% =====================================================================
%  PLANTILLA — copia este archivo como equipo_NN.tex por cada equipo
%  observado (todos menos el tuyo) y agrégalo en informe.tex con \input.
%
%  Reglas de la guía para llenar:
%  - La evidencia debe ser OBSERVABLE: n.º de diapositiva, minuto,
%    concepto explicado u omitido, error técnico, salida de la demo,
%    pregunta y respuesta. No vale "estuvo bien".
%  - Nivel 4 solo si se cumple TODO. Si dudas entre dos niveles, el inferior.
%  - Toda evidencia sale de TU hoja de campo de ese equipo (foto en hojas_campo/).
%  - 3 o más dimensiones sin evidencia = CERO.
% =====================================================================

\section{\completar{Equipo N --- Tema}}

\identificacion
  {\completar{nombre o número del equipo}}
  {\completar{T? --- nombre del tema}}
  {\completar{dd/mm/2026}}
  {\completar{hh:mm}}
  {\completar{hh:mm}}

\begin{dimensiones}
  % D1 Cumplimiento de la guía (10): conceptos clave, demo en vivo,
  %    comparación y recomendación.
  \dimension{D1 Cumplimiento de la guía (10)}
    {\faltanivel}
    {\completar{qué partes de la guía cubrieron y cuáles faltaron}}

  % D2 Fundamentación conceptual (20): nivel 4 = al menos 5 conceptos clave
  %    sin errores y acorde a la fuente oficial (RFC, OASIS, W3C, Jakarta EE...).
  %    Errores graves que bajan D2: WebSocket sobre UDP, JMS como protocolo,
  %    gRPC con JSON, TCP que conserva límites de mensaje,
  %    concurrencia = paralelismo, GET para modificar datos.
  \dimension{D2 Fundamentación conceptual (20)}
    {\faltanivel}
    {\completar{conceptos explicados (diapositiva/minuto) y errores detectados}}

  % D3 Desarrollo técnico (15): nivel 4 = demo en vivo en procesos distintos
  %    con un caso de fallo + comparación de >= 2 alternativas con >= 3
  %    criterios medibles.
  \dimension{D3 Desarrollo técnico (15)}
    {\faltanivel}
    {\completar{qué mostró la demo, si hubo caso de fallo, qué comparó}}

  % D4 Análisis y recomendación (20): nivel 4 = cuándo usar y cuándo no,
  %    >= 2 razones ligadas a un caso concreto, reconoce limitaciones.
  \dimension{D4 Análisis y recomendación (20)}
    {\faltanivel}
    {\completar{recomendación dada, razones, caso, limitaciones}}

  % D5 Comunicación (25): no lee, tiempo T +- 2 min, responde con precisión
  %    una pregunta técnica, diapositivas legibles, exponen todos.
  \dimension{D5 Comunicación (25)}
    {\faltanivel}
    {\completar{duración real, pregunta hecha y respuesta, quiénes expusieron}}

  % D6 Fuentes (10): fuente visible en cada dato/figura, fuentes primarias
  %    o indexadas, formato IEEE.
  \dimension{D6 Fuentes (10)}
    {\faltanivel}
    {\completar{diapositivas con/sin fuente, tipo de fuentes, formato}}
\end{dimensiones}

% Pon aquí los mismos niveles de la tabla, en orden D1..D6 (solo números):
\notasugerida{0}{0}{0}{0}{0}{0}

\retroalimentacion
  {\completar{una fortaleza concreta, con la evidencia que la respalda}}
  {\completar{una mejora accionable: qué hacer distinto y cómo}}

\clearpage

%% =====================================================================
%  PLANTILLA — copia este archivo como equipo_NN.tex por cada equipo
%  observado (todos menos el tuyo) y agrégalo en informe.tex con \input.
%
%  Reglas de la guía para llenar:
%  - La evidencia debe ser OBSERVABLE: n.º de diapositiva, minuto,
%    concepto explicado u omitido, error técnico, salida de la demo,
%    pregunta y respuesta. No vale "estuvo bien".
%  - Nivel 4 solo si se cumple TODO. Si dudas entre dos niveles, el inferior.
%  - Toda evidencia sale de TU hoja de campo de ese equipo (foto en hojas_campo/).
%  - 3 o más dimensiones sin evidencia = CERO.
% =====================================================================

\section{\completar{Equipo N --- Tema}}

\identificacion
  {\completar{nombre o número del equipo}}
  {\completar{T? --- nombre del tema}}
  {\completar{dd/mm/2026}}
  {\completar{hh:mm}}
  {\completar{hh:mm}}

\begin{dimensiones}
  % D1 Cumplimiento de la guía (10): conceptos clave, demo en vivo,
  %    comparación y recomendación.
  \dimension{D1 Cumplimiento de la guía (10)}
    {\faltanivel}
    {\completar{qué partes de la guía cubrieron y cuáles faltaron}}

  % D2 Fundamentación conceptual (20): nivel 4 = al menos 5 conceptos clave
  %    sin errores y acorde a la fuente oficial (RFC, OASIS, W3C, Jakarta EE...).
  %    Errores graves que bajan D2: WebSocket sobre UDP, JMS como protocolo,
  %    gRPC con JSON, TCP que conserva límites de mensaje,
  %    concurrencia = paralelismo, GET para modificar datos.
  \dimension{D2 Fundamentación conceptual (20)}
    {\faltanivel}
    {\completar{conceptos explicados (diapositiva/minuto) y errores detectados}}

  % D3 Desarrollo técnico (15): nivel 4 = demo en vivo en procesos distintos
  %    con un caso de fallo + comparación de >= 2 alternativas con >= 3
  %    criterios medibles.
  \dimension{D3 Desarrollo técnico (15)}
    {\faltanivel}
    {\completar{qué mostró la demo, si hubo caso de fallo, qué comparó}}

  % D4 Análisis y recomendación (20): nivel 4 = cuándo usar y cuándo no,
  %    >= 2 razones ligadas a un caso concreto, reconoce limitaciones.
  \dimension{D4 Análisis y recomendación (20)}
    {\faltanivel}
    {\completar{recomendación dada, razones, caso, limitaciones}}

  % D5 Comunicación (25): no lee, tiempo T +- 2 min, responde con precisión
  %    una pregunta técnica, diapositivas legibles, exponen todos.
  \dimension{D5 Comunicación (25)}
    {\faltanivel}
    {\completar{duración real, pregunta hecha y respuesta, quiénes expusieron}}

  % D6 Fuentes (10): fuente visible en cada dato/figura, fuentes primarias
  %    o indexadas, formato IEEE.
  \dimension{D6 Fuentes (10)}
    {\faltanivel}
    {\completar{diapositivas con/sin fuente, tipo de fuentes, formato}}
\end{dimensiones}

% Pon aquí los mismos niveles de la tabla, en orden D1..D6 (solo números):
\notasugerida{0}{0}{0}{0}{0}{0}

\retroalimentacion
  {\completar{una fortaleza concreta, con la evidencia que la respalda}}
  {\completar{una mejora accionable: qué hacer distinto y cómo}}

\clearpage


\end{document}
